% Section 6.8 — Additional-Source / Integrated Evaluation
% Presents EXT-3-Original and integrated SDI-4 + EXT-3-Original results.

\section{Additional-Source and Integrated Evaluation}
\label{sec:multisource}

This section evaluates the proposed classifier on the EXT-3-Original dataset (three classes, external source) and the integrated SDI-4 + EXT-3-Original benchmark. The integrated evaluation exposes the classifier to domain shift and class-composition changes, providing a more challenging generalisation test.

\subsection{EXT-3-Original Performance}
\label{sec:ext3-performance}

The EXT-3-Original dataset contains 315 images across three classes (Healthy, BG, WSSV) with a fixed partition of 221/47/47. Results are reported in Table~\ref{tab:ext3-results}.

\begin{table}[htbp]
\centering
\caption{EXT-3-Original classification results (seed 42).}
\label{tab:ext3-results}
\begin{tabular}{l C{2.2cm} C{2.2cm}}
\toprule
\textbf{Method} & \textbf{Accuracy (\%)} & \textbf{Macro-F1 (\%)} \\
\midrule
CE Baseline     & 70.21 & 72.30 \\
Proposed        & 91.49 & 91.29 \\
\midrule
$\Delta$ & +21.28 & +18.99 \\
\bottomrule
\end{tabular}
\end{table}

The proposed method achieves a substantial improvement over the CE baseline on the EXT-3-Original dataset, with a Macro-F1 of 91.29\% compared to 72.30\%. This improvement is larger than the improvement observed on SDI-4, suggesting that the ASL-LDAM + SimAM-DCFR components are particularly beneficial when the class distribution shifts or the acquisition conditions differ from the primary training source.

\subsection{Integrated SDI-4 + EXT-3-Original Performance}
\label{sec:integrated-performance}

The integrated dataset combines SDI-4 and EXT-3-Original at the partition level, yielding 1,025 training, 219 validation, and 220 test images. The test partition contains images from both sources, with the WSSV\_BG class appearing only in the SDI-4 subset. Results are reported in Table~\ref{tab:integrated-results}.

\begin{table}[htbp]
\centering
\caption{Integrated SDI-4 + EXT-3-Original classification results (seed 42).}
\label{tab:integrated-results}
\begin{tabular}{l C{2.2cm} C{2.2cm}}
\toprule
\textbf{Method} & \textbf{Accuracy (\%)} & \textbf{Macro-F1 (\%)} \\
\midrule
CE Baseline     & 87.73 & 86.51 \\
Proposed        & 83.18 & 82.18 \\
\bottomrule
\end{tabular}
\end{table}

\subsection{Discussion of the Negative Integrated Result}
\label{sec:integrated-discussion-negative}

The integrated evaluation reveals a notable negative result: the proposed method (82.18\% Macro-F1) underperforms the CE baseline (86.51\% Macro-F1) on the combined dataset. This result is a genuine observation under the recorded protocol and must be clearly acknowledged.

The source-stratified integrated accuracy provides additional insight:

\begin{itemize}
    \item Proposed EXT subset accuracy: 85.11\%
    \item Proposed SDI subset accuracy: 82.66\%
    \item CE EXT subset accuracy: 91.49\%
    \item CE SDI subset accuracy: 86.71\%
\end{itemize}

Both methods degrade on the SDI subset relative to the SDI-4-only evaluation, indicating that the domain shift introduced by EXT-3-Original affects the four-class decision boundary. The degradation is more pronounced for the proposed method, which achieved 91.01\% Macro-F1 on SDI-4 alone but only 82.66\% on the SDI subset of the integrated test.

Several hypotheses may explain this degradation, but they remain hypotheses rather than established causal explanations:

\begin{itemize}
    \item \textbf{Source interaction}: The combined training set introduces images from different acquisition conditions, potentially shifting the feature distribution learned by the backbone.
    \item \textbf{WSSV vs.~WSSV\_BG boundary}: The integrated evaluation includes only SDI-4 images for WSSV\_BG, while the WSSV class receives additional samples from EXT-3-Original. This asymmetry may alter the LDAM margin computation.
    \item \textbf{Loss dynamics}: The ASL-LDAM loss, which is sensitive to class priors, may behave differently when the prior distribution changes through the addition of EXT-3 samples.
    \item \textbf{Heterogeneous texture}: The SimAM-DCFR texture branch may overfit to source-specific textures when the training distribution becomes more heterogeneous.
\end{itemize}

This negative result does not invalidate the primary SDI-4 findings but establishes a clear limitation: the proposed method's multi-source generalisation capability is condition-dependent and requires further investigation. The CE baseline, while weaker on the individual SDI-4 and EXT-3 benchmarks, shows greater robustness to the integrated training condition.
